% 致谢页

\clearpage
\phantomsection
\addcontentsline{toc}{chapter}{致谢}

\chapter*{致谢}
时光荏苒，本论文的完成，标志着我大学学习生涯即将告一段落。
在此，我谨向所有给予我指导、帮助与支持的老师、同学和家人致以最诚挚的谢意。

感谢我的指导老师李斌，本文的完成离不开老师在选题、框架和细节上的指导。
李老师提供的因子库也是本文最重要的数据源。

感谢与我朝夕相处的同学们，感谢他们在实验中给予的帮助和建议。

感谢我的家人。他们一直是我最坚实的后盾，给予我无私的关爱、理解与支持，让我能够安心完成学业。

回望这几个月，多少次通宵达旦，对着屏幕调试、改错。
还记得最开始，每次运行后，都会有几个节点莫名其妙挂掉，
来回检查，发现是由于监控循环的间隔太短，很多节点还没来得及取数，数据就被清理了。
这件事情令我印象深刻，问题实在是很简单，在数据少的测试环境，容器加载和取数很快，
在监控触发前就取到了；但是随着容器和数据量增大，时序竞争的问题自然而然就出现了。
我会感慨，AI发展得太快了，在5年前，一个老股民可能需要花一个晚上整理公司的基本面和技术面信息，
写分析报告，制定交易策略；但现在，使用诸如AI-Hedge Fund或TradingAgent-CN之类的项目，
只需要轻轻点一下，一篇完整、详细、精准的股评报告就自动生成了，多个Agent专家相互辩论，
风格各异，结论有数据支持，多数情况可信。
在我看来，这些年的变化，就好比是从测试环境到生产环境那样。

感谢自己能生活在这个风云变幻的时代，
也感谢学校、学院和老师对我的栽培。
下一站是北京和纽约——既然前路是光明的，
那就像灯塔一样，为一切夜里不能航行的人，用火光把道路照明。

最后，感谢所有审阅本文的老师和专家，拨冗阅读拙作。
由于本人学识有限，文中难免存在疏漏与不足之处，恳请各位老师批评指正。
